\begin{textsample}{POS Dim 1 – ro – Score 59.00 – ro/ro\_ef\_51.txt}  \label{ex:f1_pos_136}
9.1.3 Língua Portuguesa O documento Curricular do Estado de Rondônia tem como principal intenção possibilitar uma língua contextualizada, que se usa efetivamente, nas práticas em que o indivíduo participa.
Para que esse \textbf{foco} no uso seja possibilitado, é preciso que a unidade linguística básica do trabalho de Língua Portuguesa seja o texto, pois é nele, materialidade do discurso, que a língua, por se encontrar em funcionamento, torna -se linguagem.
Buscando -se \textbf{aproximar} as práticas de linguagem de seus usos reais, destacamos, ainda neste documento curricular, a importância do trabalho com os \textbf{multiletramentos} fundamentais para a efetiva participação nas práticas sociais de linguagem \textbf{contemporânea}, a presença e especificidades dos textos multimodais, característicos também da cultura \textbf{digital}, e o reconhecimento da interculturalidade, \textbf{constitutiva} e \textbf{imbricada} nas diversidades pluriétnicas, tais como as seguintes comunidades : quilombolas, indígenas, ribeirinhas, surdas, ciganas, imigrantes, emigrantes, refugiadas e de deficientes visuais ( cego e baixa visão ), como \textbf{mobilização} para a superação das desigualdades de forma consubstanciada, o que impõe a necessidade de reconhecimento cultural de diferentes manifestações de conhecimento que efetivamente se comuniquem.
Para \textbf{tanto}, a BNCC ( BRASIL, 2017 ) foi elaborada com objetivo de garantir os conteúdos disciplinares comuns a serem trabalhados em todas as escolas brasileiras de Educação Básica.
No documento, idealiza -se que os \textbf{alunos} matriculados em todas as escolas de ensino básico, nas cinco regiões geográficas brasileiras, tenham acesso aos mesmos componentes curriculares ( \textbf{disciplinas} escolares ), a partir da \textbf{indicação} de objetos de conhecimento ( conteúdos disciplinares ) e de habilidades a serem trabalhadas ao longo da escolaridade.
Dessa \textbf{maneira}, a linguagem é considerada como a \textbf{capacidade} humana de articular significados coletivos e compartilhá -los em sistemas arbitrários de representação – variáveis de acordo com as múltiplas necessidades e experiências do convívio social – qualquer ato de linguagem tem como objetivo a produção de sentido.
O ensino-aprendizagem da língua materna deve, por isso, ter como objeto a produção de sentido que se dá por meio de discursos e de sua manifestação material, os textos, os quais se atualizam os códigos de uso da língua, e consequentemente, as condições de interpretação das estruturas linguísticas.
A produção de sentido por meio de textos se dá não apenas a partir dos elementos formais da língua – combinações de sons, palavras e estruturas – mas da relação dessas expressões com a situação comunicativa concreta e com o sistema cultural de representação da realidade.
É a língua, portanto, uma das formas de manifestação da linguagem, construída histórica e socialmente pelo homem.
O componente Língua Portuguesa da BNCC dialoga com documentos e orientações curriculares produzidos nas últimas décadas, buscando atualizá -los em relação às pesquisas recentes da área e às transformações das práticas de linguagem ocorridas neste \textbf{século}, devidas em grande parte ao desenvolvimento das \textbf{tecnologias} \textbf{digitais} da informação e comunicação ( TDIC ).
Assume -se aqui a perspectiva enunciativo-discursiva de linguagem, já assumida em outros documentos, como os \textbf{Parâmetros} Curriculares Nacionais ( PCN ), para os quais a linguagem é “ uma forma de ação interindividual orientada para uma finalidade específica ; um processo de interlocução que se realiza nas práticas sociais existentes numa sociedade, nos distintos momentos de sua história ” ( BRASIL, 1998, p. 20 ).
Tal proposta assume a centralidade do texto como unidade de trabalho e as perspectivas enunciativo-discursivas na \textbf{abordagem}, de forma a sempre relacionar os textos a seus contextos de produção e o desenvolvimento de habilidades ao uso significativo da linguagem em atividades de leitura, escuta e produção de textos em várias mídias e semioses.
Ao mesmo tempo que se fundamenta em concepções e conceitos já disseminados em outros documentos e orientações curriculares e em contextos variados de formação de professores, já relativamente conhecidos no ambiente escolar – tais como práticas de linguagem, discurso e gêneros discursivos/gêneros textuais, esferas/campos de circulação dos discursos –, considera as práticas \textbf{contemporâneas} de linguagem, sem o que a participação nas esferas da vida pública, do trabalho e pessoal pode se dar de forma desigual.
Na esteira do que foi proposto nos \textbf{Parâmetros} Curriculares Nacionais, o texto ganha centralidade na definição dos conteúdos, habilidades e objetivos, considerado a partir de seu pertencimento a um gênero \textbf{discursivo} que circula em diferentes esferas/campos sociais de atividade/comunicação/uso da linguagem.
Os conhecimentos sobre os gêneros, sobre os textos, sobre a língua, sobre a norma-padrão, sobre as diferentes linguagens ( semioses ) devem ser \textbf{mobilizados} em favor do desenvolvimento das \textbf{capacidades} de leitura, produção e tratamento das linguagens, que, por sua vez, devem estar a serviço da ampliação das possibilidades de participação em práticas de diferentes esferas/campos de atividades humanas.
Ao componente Língua Portuguesa cabe, então, proporcionar aos estudantes experiências que contribuam para a ampliação dos \textbf{letramentos}, de forma a possibilitar a participação significativa e \textbf{crítica} nas diversas práticas sociais permeadas/constituídas pela oralidade, pela escrita e por outras linguagens.
As práticas de linguagem \textbf{contemporâneas} não só envolvem novos gêneros e textos cada vez mais \textbf{multissemióticos} e multimidiáticos, como também novas formas de produzir, de \textbf{configurar}, de disponibilizar, de replicar e de interagir.
As novas ferramentas de \textbf{edição} de textos, áudios, fotos, \textbf{vídeos} tornam acessíveis a qualquer um a produção e disponibilização de textos \textbf{multissemióticos} nas redes sociais e outros ambientes da Web.
Não só é possível acessar conteúdos variados em diferentes mídias, como também produzir e \textbf{publicar} fotos, \textbf{vídeos} diversos, podcasts, \textbf{infográficos}, enciclopédias colaborativas, revistas e livros \textbf{digitais} etc.
Depois de ler um livro de literatura ou assistir a um filme, pode -se \textbf{postar} comentários em redes sociais específicas, seguir diretores, \textbf{autores}, escritores, acompanhar de perto seu trabalho ; podemos produzir playlists, vlogs, vídeos-minuto, escrever fanfics, produzir e-zines, nos tornar um booktuber, dentre outras muitas possibilidades.
Em \textbf{tese}, a Web é democrática : todos podem acessá -la e alimentá -la continuamente.
Mas se esse espaço é livre e bastante familiar para crianças, adolescentes e jovens de hoje, por que a escola teria que, de alguma forma, considerá -lo?
Ser familiarizado e usar não significa necessariamente levar em conta as dimensões ética, estética e política desse uso, nem tampouco lidar de forma \textbf{crítica} com os conteúdos que circulam na Web.
A contrapartida do fato de que todos podem \textbf{postar} quase tudo é que os critérios editoriais e seleção do que é adequado, bom, fidedigno não estão “ garantidos ” de início.
Passamos a \textbf{depender} de curadores ou de uma curadoria própria, que \textbf{supõe} o desenvolvimento de diferentes habilidades.
A viralização de conteúdos/publicações \textbf{fomenta} fenômenos como o da pós-verdade, em que as opiniões importam mais do que os fatos em si.
Nesse contexto, torna -se menos importante checar/verificar se algo aconteceu do que simplesmente acreditar que aconteceu ( já que isso vai ao encontro da própria opinião ou perspectiva ).
As fronteiras entre o público e o privado estão sendo recolocadas.
Não se trata de querer impor a tradição a qualquer custo, mas de refletir sobre as redefinições desses limites e de desenvolver habilidades para esse trato, inclusive refletindo sobre questões envolvendo o excesso de \textbf{exposição} nas redes sociais.
Em nome da liberdade de expressão, não se pode \textbf{dizer} qualquer coisa em qualquer situação.
Se, potencialmente, a internet seria o lugar para a divergência e o diferente circularem, na prática, a maioria das interações se dá em diferentes bolhas, em que o outro é parecido e pensa de forma semelhante.
Assim, compete à escola garantir o trato, cada vez mais necessário, com a diversidade, com a diferença.
Não se trata de deixar de privilegiar o escrito/impresso nem de deixar de considerar gêneros e práticas consagrados pela escola, tais como notícia, reportagem, \textbf{entrevista}, artigo de opinião, charge, tirinha, crônica, conto, verbete de enciclopédia, artigo de divulgação científica etc., próprios do \textbf{letramento} da letra e do impresso, mas de contemplar também os novos \textbf{letramentos}, essencialmente \textbf{digitais}.
Como resultado de um trabalho de pesquisa sobre produções culturais, é possível, por exemplo, \textbf{supor} a produção de um ensaio e de um vídeo-minuto.
No primeiro caso, um maior \textbf{aprofundamento} teórico-conceitual sobre o objeto parece necessário, e \textbf{certas} habilidades analíticas estariam mais em evidência.
No segundo caso, ainda que um nível de \textbf{análise} possa/tenha que existir, as habilidades \textbf{mobilizadas} estariam mais ligadas à síntese e percepção das potencialidades e formas de construir sentido das diferentes linguagens.
Ambas as habilidades são importantes.
Compreender uma \textbf{palestra} é importante, assim como ser capaz de atribuir diferentes sentidos a um gif ou meme.
Da mesma forma que fazer uma comunicação oral adequada e saber produzir gifs e memes significativos também podem sê -lo.
Uma parte considerável das crianças e jovens que estão na escola hoje vai exercer profissões que ainda nem existem e se deparar com \textbf{problemas} de diferentes ordens e que podem requerer diferentes habilidades, um repertório de experiências e práticas e o domínio de ferramentas que a vivência dessa diversificação pode favorecer.
O que pode parecer um gênero menor ( no sentido de ser menos valorizado, relacionado a situações tidas como pouco sérias, que envolvem paródias, chistes, remixes ou condensações e narrativas paralelas ), na verdade, pode favorecer o domínio de modos de significação nas diferentes linguagens, o que a \textbf{análise} ou produção de uma foto convencional, por exemplo, pode não propiciar.
Essa consideração dos novos e \textbf{multiletramentos} ; e das práticas da cultura \textbf{digital} no currículo não contribui somente para que uma participação mais efetiva e \textbf{crítica} nas práticas \textbf{contemporâneas} de linguagem por parte dos estudantes possa ter lugar, mas permite também que se possa ter em \textbf{mente} mais do que um “ usuário da língua/das linguagens ”, na direção do que alguns \textbf{autores} vão denominar de designer : alguém que toma algo que já existe ( inclusive textos escritos ), mescla, remixa, transforma, redistribui, produzindo novos sentidos, processo que alguns \textbf{autores} associam à criatividade.
Parte do sentido de criatividade em circulação nos dias \textbf{atuais} ( “ economias criativas ”, “ cidades criativas ” etc. ) tem algum tipo de relação com esses fenômenos de reciclagem, mistura, apropriação e redistribuição.
Dessa forma, a BNCC procura contemplar a cultura \textbf{digital}, diferentes linguagens e diferentes \textbf{letramentos}, desde aqueles basicamente lineares, com baixo nível de hipertextualidade, até aqueles que envolvem a hipermídia.
Da mesma \textbf{maneira}, \textbf{imbricada} à questão dos \textbf{multiletramentos}, essa proposta considera, como uma de suas premissas, a diversidade cultural.
Sem aderir a um \textbf{raciocínio} classificatório reducionista, que desconsidera as hibridizações, apropriações e mesclas, é importante contemplar o cânone, o marginal, o culto, o popular, a cultura de massa, a cultura das mídias, a cultura \textbf{digital}, as culturas infantis e juvenis, de forma a garantir uma ampliação de repertório e uma interação e trato com o diferente.
Ainda em relação à diversidade cultural, cabe \textbf{dizer} que se estima que mais de 250 línguas são faladas no país – indígenas, de imigração, de sinais, crioulas e afro-brasileiras, além do português e de suas variedades.
Esse patrimônio cultural e linguístico é desconhecido por grande parte da população brasileira.
No Brasil com a Lei no 10.436, de 24 de abril de 2002, oficializou -se também a Língua Brasileira de Sinais ( \textbf{Libras} ), tornando possível, em âmbito nacional, realizar discussões relacionadas à necessidade do respeito às particularidades linguísticas da comunidade surda e do uso dessa língua nos ambientes escolares.
Assim, é \textbf{relevante} no espaço escolar conhecer e valorizar as realidades nacionais e internacionais da diversidade linguística e analisar diferentes situações e atitudes humanas \textbf{implicadas} nos usos linguísticos, como o preconceito linguístico.
Por outro \textbf{lado}, existem muitas línguas ameaçadas de extinção no país e no mundo, o que nos chama a atenção para a \textbf{correlação} entre repertórios culturais e linguísticos, pois o desaparecimento de uma língua impacta significativamente a cultura.
Muitos representantes de comunidades de falantes de diferentes línguas, especialistas e \textbf{pesquisadores} vêm \textbf{demandando} o reconhecimento de direitos linguísticos.
Por isso, já temos municípios brasileiros que cooficializaram línguas indígenas – tukano, baniwa, nheengatu, akwexerente, guarani, macuxi – e línguas de migração – talian, pomerano, hunsrickisch –, existem publicações e outras ações expressas nessas línguas ( livros, jornais, filmes, peças de teatro, programas de radiodifusão ) e programas de educação bilíngue.
As recomendações para o componente curricular de Língua Portuguesa objetivam “ garantir a todos os \textbf{alunos} o acesso aos saberes linguísticos necessários para a participação social e o exercício da cidadania ” ( BRASIL, 2017, p. 63 ).
A linguagem é o pressuposto necessário para que as pessoas sejam capazes de elaborar pensamentos, comunicar -se, acessar os mais diversos tipos de informações e produzir conhecimentos.
Considerando o texto como “ centro das práticas de linguagem ”, os objetos de conhecimento/conteúdo para o ensino de Língua Portuguesa estão organizados em quatro eixos \textbf{organizadores} : oralidade, leitura/escuta, produção ( escrita e \textbf{multissemiótica} ) e \textbf{análise} linguística/semiótica ( que envolve conhecimentos linguísticos – sobre o sistema de escrita, o sistema da língua e a norma-padrão –, textuais, \textbf{discursivos} e sobre os modos de organização e os elementos de outras semioses ) ( BNCC, 2017 ).
FONTE : Elaborado pela Equipe Redatora de Língua Portuguesa.
Cada eixo está organizado em unidades temáticas, que, por sua vez, estão atrelados a objetos de conhecimento e habilidades.
Cabe ressaltar, \textbf{reiterando} o movimento \textbf{metodológico} de documentos curriculares anteriores, que estudos de natureza \textbf{teórica} e metalinguística – sobre a língua, sobre a literatura, sobre a norma padrão e outras variedades da língua – não devem nesse nível de ensino ser tomados como um fim em si mesmo, devendo estar envolvidos em práticas de reflexão que permitam aos estudantes ampliarem suas \textbf{capacidades} de uso da língua/linguagens ( em leitura e em produção ) em práticas situadas de linguagem.
O Eixo Leitura compreende as práticas de linguagem que \textbf{decorrem} da interação ativa do leitor/ouvinte/espectador com os textos escritos, orais e \textbf{multissemióticos} e de sua interpretação, sendo exemplos as leituras para : fruição estética de textos e obras literárias ; pesquisa e embasamento de trabalhos escolares e acadêmicos ; realização de procedimentos ; conhecimento, discussão e \textbf{debate} sobre temas sociais \textbf{relevantes} ; sustentar a reivindicação de algo no contexto de atuação da vida pública ; ter mais conhecimento que permita o desenvolvimento de projetos pessoais, dentre outras possibilidades.
Leitura no contexto da BNCC é tomada em um sentido mais amplo, \textbf{dizendo} respeito não somente ao texto escrito, mas também a imagens estáticas ( foto, pintura, desenho, esquema, gráfico, diagrama ) ou em movimento ( filmes, \textbf{vídeos} etc. ) e ao som ( música ), que acompanha e cossignifica em muitos gêneros \textbf{digitais}.
O tratamento das práticas leitoras compreende dimensões \textbf{inter-relacionadas} às práticas de uso e reflexão, tais como as apresentadas a seguir ( BNCC, 2017 ).

% matched lemmas: abordagem, aluno, análise, aprofundamento, aproximar, atual, autor, capacidade, certo, configurar, constitutivo, contemporâneo, correlação, crítico, debate, decorrer, demandar, depender, digital, disciplina, discursivo, dizer, edição, entrevista, exposição, foco, fomentar, imbricar, implicar, indicação, infográfico, inter-relacionar, lado, letramento, libra, maneira, mente, metodológico, mobilizar, mobilização, multiletramento, multissemiótico, organizador, palestra, parâmetro, pesquisador, postar, problema, publicar, raciocínio, reiterar, relevante, supor, século, tanto, tecnologia, tese, teórico, vídeo
\end{textsample}
